\documentclass[11pt]{article}
\usepackage{amsmath,amssymb,amsthm,mathtools}
\usepackage[margin=1in]{geometry}
\usepackage{longtable}
\usepackage{url}

\newtheorem{theorem}{Theorem}
\newtheorem{lemma}{Lemma}
\newcommand{\dd}{\,d}
\newcommand{\GammaH}{\Gamma_H}

\title{Arbitrary-Graphon Stability at the High-Density Endpoint}
\author{Exact certificate note for the remaining Atlas rows}
\date{August 14, 2026}

\begin{document}
\maketitle

\section{Result and relation to the two endpoint regimes}

Let $H$ be a finite simple graph with $n$ vertices and $m$ edges.  For a
graphon $W$ put
\[
 p(W)=\int_{\Omega^2}W,
 \qquad
 \Phi_H(p)=(1-p)^n\chi_H\!\left(\frac1{1-p}\right),
 \qquad
 \GammaH(W)=t(H,W)-\Phi_H(p(W)).
\]
Write $V=1-W$ and $q=\int V=1-p(W)$.  The preceding maximum-degree theorem
proves $\GammaH(1-V)\geq0$ when the complement degree is uniformly small.
The preceding complete-core theorem proves the same conclusion when every
defect is incident to one small exceptional set.  The purpose of this note is
to combine those regimes by applying the decomposition to the high-degree set
of an arbitrary complement.

The outcome is a universal endpoint theorem on the 29-row pre-Atlas-152
audit snapshot: for each row there is an explicit rational
$\varepsilon_H>0$ such that
\[
 p(W)\geq1-\varepsilon_H
 \quad\Longrightarrow\quad
 t(H,W)\geq\Phi_H(p(W))
\]
for every graphon on every probability space.  There is no regularity,
finite-step, pointwise kernel, or exact complete-core assumption.  The radii
are deliberately conservative; their role is to prove a genuine uniform
neighbourhood, not to estimate the optimal endpoint.
Atlas $152$ was subsequently classified negative near the threshold; the
other $28$ rows are the current open set and inherit this endpoint theorem.

\section{Constants from the two component theorems}

Expand the chromatic target at the endpoint as
\begin{equation}
 \Phi_H(1-q)=1-mq+R_H(q),
 \qquad
 R_H(q)=\sum_{j=2}^{n}b_jq^j.
 \label{eq:target}
\end{equation}
Let
\[
 a_H=2m-n,
 \qquad
 g_H=\sum_{x\in V(H)}\binom{\deg_H(x)}2-N_3(H),
\]
and let $\Lambda_H,\Omega_H$ be the exact higher-rank constants from
\path{notes/maximum_degree_complement_high_density_stability.tex}.  Thus, if
$F$ is a graphon with complement density $u$, maximum complement degree
$\Delta_F$, and
\[
 X_F=t(P_3,F)-u^2,
 \qquad
 Y_F=t(P_3,F)-t(K_3,F),
\]
then
\begin{equation}
 \GammaH(1-F)\geq
 (g_H-\Lambda_H\Delta_F)X_F
 +(N_3(H)-\Omega_H\Delta_F)Y_F.
 \label{eq:max-degree}
\end{equation}

Define also
\begin{align}
 D_H&=\sum_{j=2}^{n}j|b_j|,\\
 J_H&=\sum_{j=2}^{n}j(j-1)|b_j|,\label{eq:J}\\
 \eta_H&=\min\left\{\frac{g_H}{\Lambda_H},
                       \frac{N_3(H)}{\Omega_H}\right\},
 \qquad
 \delta_H=\frac{\eta_H}{2},\label{eq:delta}\\
 \tau_H&=\min\left\{
 \frac12,
 \frac{g_H}{2\Lambda_H},
 \frac{N_3(H)}{\Omega_H},
 \frac{a_Hg_H}
 {a_Hg_H(n-1)+4D_Hg_H+8D_H^2}
 \right\}.\label{eq:tau}
\end{align}
All denominators in these formulas are positive for the $29$ rows considered
below.

\section{Two elementary interaction lemmas}

\begin{lemma}[Mixed chromatic-target remainder]
\label{lem:target}
If $x,y\geq0$ and $x+y\leq1$, then
\[
 |R_H(x+y)-R_H(x)-R_H(y)|\leq J_Hxy.
\]
\end{lemma}

\begin{proof}
For every $j\geq2$, two integrations of the mixed derivative give
\[
 (x+y)^j-x^j-y^j
 =j(j-1)\int_0^x\int_0^y(u+v)^{j-2}\dd v\dd u.
\]
The integrand is at most one because $u+v\leq x+y\leq1$.  Multiply by
$|b_j|$ and sum.
\end{proof}

For a finite collection $\mathcal E$ of numbers in $[0,1]$, put
\[
 \mathcal F(\mathcal E)=
 \prod_{z\in\mathcal E}(1-z)-1+\sum_{z\in\mathcal E}z.
\]

\begin{lemma}[Superadditive union remainder]
\label{lem:union}
For disjoint collections $\mathcal E_0,\mathcal E_1$,
\[
 \mathcal F(\mathcal E_0\cup\mathcal E_1)
 \geq\mathcal F(\mathcal E_0)+\mathcal F(\mathcal E_1).
\]
Consequently, if graphons $V_0,V_1$ have disjoint pointwise supports and
$q_i=\int V_i$, then
\begin{equation}
 \GammaH(1-V_0-V_1)
 \geq\GammaH(1-V_0)+\GammaH(1-V_1)
 -J_Hq_0q_1.
 \label{eq:split}
\end{equation}
\end{lemma}

\begin{proof}
If $P_i=\prod_{z\in\mathcal E_i}(1-z)$, direct expansion gives
\[
 \mathcal F(\mathcal E_0\cup\mathcal E_1)
 =\mathcal F(\mathcal E_0)+\mathcal F(\mathcal E_1)
 +(1-P_0)(1-P_1).
\]
Apply this identity pointwise to the edge factors in $t(H,1-V_0-V_1)$ and
integrate.  The linear terms cancel.  Subtracting the target leaves
$R_H(q_0)+R_H(q_1)-R_H(q_0+q_1)$, which is bounded below by
Lemma~\ref{lem:target}.
\end{proof}

\section{A quantitative margin for a complete core}

Suppose $\Omega=A\sqcup B$, $\mu(A)=s$, and $V_A$ is supported on pairs with
at least one endpoint in $A$.  Thus $1-V_A=1$ on $B^2$.  Normalize the
restrictions of the measure and define
\[
 a(z)=\int_BV_A(z,y)\dd\mu_B(y),
 \qquad
 E_2=\int_Aa(z)^2\dd\mu_A(z).
\]

\begin{lemma}[Quantitative complete-core margin]
\label{lem:hub-margin}
Assume that $H$ has no isolated vertices and $a_H,g_H,N_3(H)>0$.
If $0<s\leq\tau_H/2$, then
\begin{equation}
 \GammaH(1-V_A)\geq\frac{a_H}{2}sE_2.
 \label{eq:hub-margin}
\end{equation}
\end{lemma}

\begin{proof}
The proof of the complete-core theorem, reproduced at the level needed here,
splits $V_A$ into its $A^2$ and cut parts.  If
\[
 \bar r=\int_{A^2}V_A\dd\mu_A^2,
\]
the maximum-degree theorem on the $A^2$ part, the exact
one-exceptional-vertex cut gain, and Young's inequality give
\begin{align*}
 \GammaH(1-V_A)
 &\geq \frac{g_H}{4}s^3(1-s)\bar r^2\\
 &\quad+sE_2\left[
 a_H(1-s)^{n-1}-4D_Hs
 -\frac{4D_H^2s^2}{g_H(1-s)}
 \right].
\end{align*}
The first three restrictions in \eqref{eq:tau} justify the maximum-degree
input.  For $s\leq1/2$, Bernoulli's inequality and
$(1-s)^{-1}\leq2$ bound the bracket below by
\[
 a_H-\left(a_H(n-1)+4D_H+\frac{8D_H^2}{g_H}\right)s.
\]
The final restriction in \eqref{eq:tau}, together with
$s\leq\tau_H/2$, makes this at least $a_H/2$.  Discard the nonnegative
$\bar r^2$ term.
\end{proof}

\section{Universal arbitrary-graphon endpoint theorem}

Define
\begin{equation}
 \varepsilon_H=\min\left\{
 \frac{\delta_H^2}{2},
 \frac{\delta_H\tau_H}{2},
 \frac{a_H\delta_H^2}{16J_H}
 \right\}.
 \label{eq:epsilon}
\end{equation}

\begin{theorem}[Universal endpoint]
\label{thm:universal}
Suppose
\[
 a_H,g_H,N_3(H),\Lambda_H,\Omega_H,J_H>0.
\]
Assume also that $H$ has no isolated vertices.  For every graphon $W$,
\[
 1-p(W)\leq\varepsilon_H
 \quad\Longrightarrow\quad
 t(H,W)\geq\Phi_H(p(W)).
\]
\end{theorem}

\begin{proof}
Put $V=1-W$, $q=\int V$, and $\delta=\delta_H$.  The case $q=0$ is
immediate.  Otherwise let
\[
 A=\{x:d_V(x)>\delta\},
 \qquad B=\Omega\setminus A,
 \qquad s=\mu(A).
\]
Markov's inequality gives
\begin{equation}
 s\leq\frac q\delta.
 \label{eq:markov}
\end{equation}
Split
\[
 V_A=V\mathbf1_{(A\times\Omega)\cup(\Omega\times A)},
 \qquad
 V_B=V\mathbf1_{B^2},
 \qquad q_A=\int V_A,\quad q_B=\int V_B.
\]
These kernels have disjoint supports and $q_A+q_B=q$.
If $s=0$, then $V_A=0$ almost everywhere and the maximum-degree estimate
applied to $V_B=V$ proves the result.  Hence assume $s>0$.

The first two restrictions in \eqref{eq:epsilon} and \eqref{eq:markov} give
\[
 s\leq\frac\delta2,
 \qquad
 s\leq\frac{\tau_H}{2}.
\]
For almost every $z\in A$, write
\[
 d_V(z)=s r(z)+(1-s)a(z),
 \qquad 0\leq r(z)\leq1.
\]
Since $d_V(z)>\delta$ and $s\leq\delta/2$,
\[
 a(z)>\frac{\delta-s}{1-s}\geq\frac\delta2.
\]
Consequently $E_2\geq\delta^2/4$, and
Lemma~\ref{lem:hub-margin} yields
\begin{equation}
 \GammaH(1-V_A)\geq\frac{a_H\delta^2}{8}s.
 \label{eq:A-margin}
\end{equation}

On $B$, the degree of $V_B$ is at most the degree of $V$, hence at most
$\delta=\eta_H/2$.  Equation~\eqref{eq:max-degree} therefore gives
\begin{equation}
 \GammaH(1-V_B)\geq0.
 \label{eq:B-gap}
\end{equation}
Also $q_A\leq2s$ and $q_B\leq q$.  Lemma~\ref{lem:union},
\eqref{eq:A-margin}, and \eqref{eq:B-gap} now imply
\[
 \GammaH(W)
 \geq \frac{a_H\delta^2}{8}s-J_Hq_Aq_B
 \geq \left(\frac{a_H\delta^2}{8}-2J_Hq\right)s.
\]
The final restriction in \eqref{eq:epsilon} makes the last expression
nonnegative.
\end{proof}

\section{Exact certificates for the 29-row audit snapshot}

The following table records the four new quantities.  The graph identities,
edge lists, chromatic polynomials, and the constants
$g_H,\Lambda_H,N_3(H),\Omega_H$ are independently recorded in
\path{notes/maximum_degree_complement_high_density_stability.tex}; the exact
$\tau_H$ reconstruction is recorded in
\path{notes/arbitrary_fibre_hub_high_density_stability.tex}.  Thus the table
does not depend on a drawing or a numerical graph match.

\small
\begin{longtable}{r@{\quad}r@{\quad}c@{\quad}c@{\quad}c}
Atlas&$J_H$&$\delta_H$&$\tau_H$&$\varepsilon_H$\\ \hline
43&190&$4/161$&$7/9825$&$1/703570$\\
118&586&$11/962$&$22/118923$&$121/1084620368$\\
122&586&$5/461$&$5/29668$&$25/249074612$\\
124&586&$5/464$&$5/29668$&$25/252326912$\\
126&586&$5/454$&$5/29668$&$25/241567952$\\
127&668&$9/908$&$9/71998$&$81/1101483904$\\
130&460&$1/94$&$8/39637$&$1/8129120$\\
137&774&$7/1173$&$35/204311$&$245/8519752368$\\
139&774&$13/2268$&$13/81593$&$845/31850558208$\\
145&982&$1/173$&$1/8629$&$5/235122224$\\
147&982&$13/2368$&$13/120647$&$845/44051922944$\\
148&982&$13/2414$&$13/120647$&$845/45780022976$\\
151&1190&$1/276$&$65/836593$&$1/145039104$\\
152&982&$6/1159$&$3/30122$&$45/2638203884$\\
153&1064&$3/587$&$3/34127$&$45/2932971328$\\
157&1088&$17/5450$&$17/128700$&$51/7603840000$\\
160&1088&$4/1319$&$8/64275$&$3/473214992$\\
168&1504&$17/5786$&$17/222218$&$867/201402420736$\\
169&1378&$4/1389$&$8/94731$&$2/443100723$\\
171&1586&$8/2903$&$2/30209$&$24/6682935337$\\
174&1794&$1/436$&$8/150115$&$1/454709632$\\
178&1528&$5/3069$&$35/363759$&$175/115134934464$\\
181&1818&$5/3241$&$35/485959$&$25/21824479152$\\
185&2108&$5/3299$&$35/625831$&$175/183537674464$\\
188&2316&$2/1355$&$7/147950$&$7/8504467800$\\
194&2838&$3/3662$&$48/1100605$&$3/12686090824$\\
196&2838&$3/3713$&$48/1100605$&$3/13041905074$\\
199&3128&$3/3748$&$16/435551$&$9/43940592512$\\
203&3858&$7/16158$&$63/1963067$&$49/895333652544$\\
\end{longtable}
\normalsize

Every entry is positive.  The radii range from
$49/895333652544$ to $1/703570$.  Therefore no arbitrary graphon sequence for
any of these rows can be a counterexample while $p\to1$.

\section{Reproducibility and formalization boundary}

Run
\begin{verbatim}
python codes/principled_negative_tests.py
\end{verbatim}
from the repository root.  The line headed ``arbitrary-graphon universal
high-density theorem'' reconstructs $J_H,\delta_H,\tau_H,\varepsilon_H$ from
the exact chromatic polynomial and the two independently audited component
tables.  It asserts exact equality with every row above.

A formalization can be separated into the following reusable layers.
\begin{enumerate}
 \item Prove Lemma~\ref{lem:target} for a finite polynomial and
 Lemma~\ref{lem:union} for finite products.
 \item Import or restate the maximum-degree inequality
 \eqref{eq:max-degree} and the quantitative complete-core margin
 \eqref{eq:hub-margin}.
 \item Define the measurable high-degree set $A$ and use
 $\mu(A)\leq q/\delta$.
 \item Verify the elementary bounds $E_2\geq\delta^2/4$,
 $q_A\leq2s$, and $q_B\leq q$.
 \item Reduce the analytic theorem to the rational inequalities in
 \eqref{eq:epsilon}; certify the finite table separately.
\end{enumerate}
All integrations involve bounded nonnegative kernels, so Tonelli applies
directly.  No approximation, compactness, or limiting passage is used.

This endpoint theorem alone does not settle any full row: intermediate edge
densities remain open.  Accordingly, no catalogue status changes.

\end{document}
